1|\documentclass[11pt,a4paper]{article}
2|
3|% === Layout ===
4|\usepackage[margin=2.5cm]{geometry}
5|\usepackage{amsmath,amssymb,amsthm}
6|\usepackage{mathtools}
7|\usepackage{bm}
8|\usepackage{graphicx}
9|\usepackage{hyperref}
10|\usepackage{cleveref}
11|\usepackage{booktabs}
12|\usepackage{enumitem}
13|\usepackage{datetime}
14|\usepackage{xcolor}
15|\usepackage{mathrsfs}
16|
17|% === Theorem environments ===
18|\newtheorem{theorem}{Theorem}[section]
19|\newtheorem{lemma}[theorem]{Lemma}
20|\newtheorem{corollary}[theorem]{Corollary}
21|\newtheorem{proposition}[theorem]{Proposition}
22|\newtheorem{conjecture}[theorem]{Conjecture}
23|\theoremstyle{definition}
24|\newtheorem{definition}[theorem]{Definition}
25|\newtheorem{example}[theorem]{Example}
26|\theoremstyle{remark}
27|\newtheorem{remark}[theorem]{Remark}
28|
29|% === Macros ===
30|\newcommand{\R}{\mathbb{R}}
31|\newcommand{\N}{\mathbb{N}}
32|\newcommand{\E}{\mathbb{E}}
33|\newcommand{\Var}{\mathrm{Var}}
34|\newcommand{\Cov}{\mathrm{Cov}}
35|\newcommand{\TV}{\mathrm{TV}}
36|\newcommand{\norm}[1]{\left\lVert#1\right\rVert}
37|\newcommand{\inner}[2]{\langle #1, #2 \rangle}
38|\newcommand{\abs}[1]{\left\lvert#1\right\rvert}
39|\newcommand{\dd}{\,\mathrm{d}}
40|\newcommand{\D}{\mathcal{D}}
41|\newcommand{\F}{\mathcal{F}}
42|\newcommand{\Osc}{\mathcal{O}}
43|\newcommand{\Obs}{\mathcal{O}_{\text{obs}}}
44|\newcommand{\Res}{\mathcal{R}}
45|\newcommand{\Turb}{\mathcal{T}}
46|\newcommand{\uN}{\bm{u}_N}
47|\newcommand{\uNS}{\bm{u}_{\text{NS}}}
48|\newcommand{\uRANS}{\bm{u}_{\text{RANS}}}
49|\newcommand{\uTrue}{\bm{u}_{\text{true}}}
50|\newcommand{\epsRes}{\varepsilon_{\text{res}}}
51|\newcommand{\epsStat}{\varepsilon_{\text{stat}}}
52|\newcommand{\epsModel}{\varepsilon_{\text{model}}}
53|\newcommand{\DeltaX}{\Delta x}
54|\newcommand{\Kolmogorov}{\eta}
55|\newcommand{\Reynolds}{\mathit{Re}}
56|\newcommand{\ind}{\mathbf{1}}
57|\newcommand{\rigorFull}{\textsc{[Rigor: Full]}}
58|
59|\title{\bfseries Turbulence as an Unidentifiability Theorem:\\
60|Indistinguishability of N-Body Truncation Error and True Turbulence at Finite Resolution\\
61|\normalsize Turbulence as an Unidentifiability Theorem: \\
62|Indistinguishability of Truncation Error and True Turbulence \\
63|under Finite-Resolution N-Body Dynamics}
64|
65|\author{SCX Theory Group}
66|\date{\today}
67|
68|\begin{document}
69|\maketitle
70|
71|\begin{abstract}
We propose a fundamental epistemological reconstruction: turbulence is not a closure problem, but an \textbf{unidentifiability problem}（unidentifiability problem）。
The starting point is extremely simple — Newton's second law $\bm{F}=m\bm{a}$ rigorously describes the motion of every single molecule. $N$-body molecular dynamics (MD) should converge to the Boltzmann equation in the thermodynamic limit $N\to\infty$, and subsequently to the Navier--Stokes (NS) equations. But no computational device in the universe can simulate $10^{23}$ molecules. Any finite-resolution simulation necessarily produces truncation error.

This paper presents three theorems. \textbf{Theorem~1} (Unidentifiability Theorem): Under finite-resolution observation, true turbulence, truncation error artifacts, and over-smoothed RANS fields are indistinguishable in total variation distance — $\TV(P_h^{\mathcal{W}_A}, P_h^{\mathcal{W}_B}) \leq C\cdot h^{\alpha}$, with $\TV\to 0$ as $h\to 0$, but $\TV>0$ for any finite $h>0$. \textbf{Theorem~2} (Multi-Model Audit Theorem): The probability that $M$ independent turbulence models (RANS $k$-$\varepsilon$, $k$-$\omega$ SST, LES, DES, DNS coarse grid) simultaneously produce systematic bias exceeding $\Delta$ satisfies $\mathbb{P}(\cap_{m=1}^{M}\{B_m > \Delta\}) \leq \exp(-2M\Delta^2/L^2)$, with DNS fine grid as the truth anchor — proving that model consensus reliability grows exponentially with model count. \textbf{Theorem~3} (Kolmogorov Spectral Truncation Theorem): At finite resolution $h$, the observed energy spectrum strictly decomposes as $E_h(k) = E_{\text{true}}(k) + C_{\text{res}}\cdot h^{2/3}\cdot k^{-5/3} + C_{\text{stat}}\cdot N^{-1/2} + o(h^{2/3} + N^{-1/2})$. When the truncation error term and statistical fluctuation term have comparable magnitude in the inertial range, the $-5/3$ spectrum emerges from truncation error — the Kolmogorov spectrum is not a physical law, but a mathematical inevitability at finite resolution.

The corollary: turbulence modeling is not about better closure models — it is about declaring your \textbf{resolution assumption}. At different resolutions, turbulence is different things.This structure is isomorphic to the quantum measurement problem: the observation scale determines what you see.
78|\end{abstract}
79|
80|% =================================================================
81|\section{Starting Point: A Forgotten Fact}
82|% =================================================================
83|
84|Every fluid dynamicist learns the Navier--Stokes equations:
85|\begin{equation}\label{eq:ns}
86|  \frac{\partial \bm{u}}{\partial t} + (\bm{u}\cdot\nabla)\bm{u}
87|  = -\frac{1}{\rho}\nabla p + \nu \nabla^2 \bm{u} + \bm{f},
88|\end{equation}
89|and that these emerge from Boltzmann via Chapman--Enskog. Boltzmann itself arises from the BBGKY hierarchy of the $N$-body Liouville equation. The chain bottom is Newton:
90|\begin{equation}\label{eq:newton}
91|  m_i \frac{\dd^2 \bm{x}_i}{\dd t^2} = \bm{F}_i
92|  = -\sum_{j\neq i} \nabla_{\bm{x}_i} V(\abs{\bm{x}_i - \bm{x}_j}),
93|\end{equation}
94|其中 $V(r)$ 是分子间势（例如 Lennard-Jones 势），$i=1,\dots,N$，$N\sim 10^{23}$。
95|
96|Key: Equation~\eqref{eq:newton} is \textbf{精确的}。exact within classical mechanics. If you solved $10^{23}$ coupled ODEs exactly, you would obtain \textbf{the complete flow field at all scales} — including all fluctuations we call ``turbulence.''
97|
98|But you cannot.
99|
100|This is the root. Physics is not missing anything — it is \textbf{the fundamental limit of computation}preventing a first-principles solution. This is not engineering — it is \textbf{epistemological}。
101|
102|% =================================================================
103|\section{From N-Body MD to Continuum: The Convergence Chain}
104|% =================================================================
105|
106|\subsection{Micro → Meso → Macro}
107|
108|Define the microscopic phase-space density:
109|\begin{equation}\label{eq:liouville}
110|  f_N(\bm{x}_1,\dots,\bm{x}_N,\bm{p}_1,\dots,\bm{p}_N,t),
111|\end{equation}
112|satisfying $\partial_t f_N = \{H, f_N\}$. Integrating over $N-k$ coordinates:
113|\begin{equation}\label{eq:bggky}
114|  f^{(k)}(\bm{x}_1,\dots,\bm{x}_k,\bm{p}_1,\dots,\bm{p}_k,t)
115|  = \int f_N \,\dd\bm{x}_{k+1}\dots\dd\bm{x}_N\dd\bm{p}_{k+1}\dots\dd\bm{p}_N.
116|\end{equation}
117|
118|This yields the BBGKY hierarchy. In the Boltzmann--Grad limit ($N\to\infty$, $d\to 0$, $Nd^2 = \text{const}$, $d$ = effective diameter), $f^{(1)}$ satisfies Boltzmann.然后通过 Chapman--Enskog 展开，取零阶得到 Euler 方程，取一阶得到 Navier--Stokes 方程。
119|
120|\textbf{Convergence Proposition:}
121|\begin{equation}\label{eq:convergence}
122|  \lim_{N\to\infty} \uN(t,\bm{x}) = \uNS(t,\bm{x}),
123|\end{equation}
124|where $\uN$ is the coarse-grained velocity from $N$-body MD (averaged over $\Delta x$, $\Delta t$), $\uNS$ is a NS solution.
125|
126|\subsection{The Coarse-Graining Operator}
127|
128|定义The Coarse-Graining Operator $\mathcal{C}_{h}$：
129|\begin{equation}\label{eq:coarse}
130|  (\mathcal{C}_{h}\bm{v})(t,\bm{x})
131|  = \frac{1}{\abs{B_h(\bm{x})}}
132|   \int_{B_h(\bm{x})}
133|   \frac{1}{N_h(\bm{y})}\sum_{i: \bm{x}_i(t)\in B_h(\bm{y})} \dot{\bm{x}}_i(t)
134|   \,\dd\bm{y},
135|\end{equation}
136|其中 $B_h(\bm{x})$ 是以 $\bm{x}$ 为中心、尺度 $h$ 的粗粒化核，$N_h(\bm{y})$ 是 $B_h(\bm{y})$ 内的粒子数。
137|Then:
138|\begin{equation}\label{eq:un_def}
139|  \uN(t,\bm{x}) = (\mathcal{C}_{h}\dot{\bm{X}})(t,\bm{x}),
140|\end{equation}
141|where $\dot{\bm{X}} = (\dot{\bm{x}}_1,\dots,\dot{\bm{x}}_N)$ is the $N$-body phase-space trajectory.
142|
143|% =================================================================
144|\section{Error Decomposition and Scaling Analysis at Finite Resolution}
145|% =================================================================
146|
147|This section promotes scaling analysis from appendix to main text.
148|
149|\subsection{Definition of Three Velocity Fields}
150|
151|We define three \textbf{conceptually distinct} velocity fields:
152|
153|\begin{enumerate}[label=(\arabic*),leftmargin=*]
154|  \item $\uN(t,\bm{x})$ —— $N$ 体 MD 在有限 $N$ 和有限粗粒化尺度 $h$ 下的粗粒化速度场。
155|  \item $\uNS(t,\bm{x})$ —— Navier--Stokes 方程的精确解（假设其存在且唯一，至少在统计意义上）。
156|  \item $\uRANS(t,\bm{x})$ —— 某个 RANS 湍流模型预测的（系综平均）速度场。
157|\end{enumerate}
158|
159|Note: $\uNS$ \textbf{already contains turbulence} — if NS indeed describes turbulence. $\uRANS$ is the \textbf{smoothed field} after closure — its turbulence has been modeled away.
160|
161|\subsection{Error Decomposition}
162|
163|For finite $N$ and finite $h$:
164|\begin{equation}\label{eq:error_decomp}
165|  \norm{\uN - \uNS}_{L^2}
166|  = \norm{(\uN - \E[\uN]) + (\E[\uN] - \uNS)}_{L^2}
167|  \leq \norm{\uN - \E[\uN]}_{L^2} + \norm{\E[\uN] - \uNS}_{L^2}.
168|\end{equation}
169|
170|Define:
171|\begin{align}
172|  \epsStat &\coloneqq \norm{\uN - \E[\uN]}_{L^2}
173|            = \text{统计涨落（statistical fluctuation）}, \label{eq:eps_stat}\\
174|  \epsRes  &\coloneqq \norm{\E[\uN] - \uNS}_{L^2}
175|            = \text{分辨率截断误差（resolution truncation error）}. \label{eq:eps_res}
176|\end{align}
177|
178|\subsection{Spectral Representation and Scaling of Truncation Error}
179|
180|设 $\uNS$ 的空间 Fourier 变换为 $\hat{\bm{u}}(k)$。The Coarse-Graining Operator $\mathcal{C}_h$ 在谱域等价于与滤波器 $\hat{G}_h(k)$ 的乘积：
181|\begin{equation}\label{eq:filter_k}
182|  \widehat{\mathcal{C}_h \uNS}(k) = \hat{G}_h(k) \hat{\bm{u}}(k),
183|\end{equation}
184|其中 $\hat{G}_h(0)=1$ 且 $\hat{G}_h(k) \to 0$ 当 $k \gg 2\pi/h$。常见滤波器包括 Gaussian 滤波器 $\hat{G}_h(k) = e^{-(kh)^2/2}$ 和盒式滤波器 $\hat{G}_h(k) = \operatorname{sinc}(kh/2)$。
185|
186|The $L^2$ norm of truncation error is, by Parseval:
187|\begin{equation}\label{eq:eps_res_spectral}
188|  \epsRes^2 = \norm{\E[\uN] - \uNS}_{L^2}^2
189|  = \int_{\R^3} \abs{1 - \hat{G}_h(k)}^2 \, E(k)\,\dd k,
190|\end{equation}
191|where $E(k)$ is the energy spectrum. For Kolmogorov $E(k) = C_K \epsilon^{2/3} k^{-5/3}$, Gaussian filter, $x = kh$:
192|\begin{equation}\label{eq:eps_res_kolmogorov}
193|  \epsRes^2 \sim C_K \epsilon^{2/3} \int_{0}^{\infty} (1 - e^{-(kh)^2/2})^2 k^{-5/3}\,\dd k
194|  = C_K \epsilon^{2/3} h^{2/3} \int_{0}^{\infty} (1 - e^{-x^2/2})^2 x^{-5/3}\,\dd x.
195|\end{equation}
196|Since $\int_{0}^{\infty} (1 - e^{-x^2/2})^2 x^{-5/3}\,\dd x < \infty$:
197|\begin{equation}\label{eq:eps_res_scaling_final}
198|  \epsRes \sim C_{\text{res}} \cdot h^{1/3},
199|  \quad C_{\text{res}} = \sqrt{C_K}\,\epsilon^{1/3}\left(\int_{0}^{\infty} (1 - e^{-x^2/2})^2 x^{-5/3}\,\dd x\right)^{1/2}.
200|\end{equation}
201|Generally, if $E(k) \sim k^{-\beta}$, then $\epsRes \sim h^{(\beta-1)/2}$.
202|
203|\subsection{Spectral Representation and Scaling of Statistical Fluctuations}
204|
205|Statistical fluctuations within $V_h = h^3$ from finite molecular sampling. Maxwell--Boltzmann with $\sigma_v^2 = k_B T/m$. $N_h \sim n h^3$, $n = N/V$. By i.i.d.\ propagation:
206|\begin{equation}\label{eq:eps_stat_detail}
207|  \epsStat^2 \sim \frac{\sigma_v^2}{N_h} \sim \frac{k_B T/m}{n h^3}
208|  \sim \frac{k_B T}{\rho} \cdot \frac{1}{h^3},
209|\end{equation}
210|where $\rho = mn$ is mass density. Thus:
211|\begin{equation}\label{eq:eps_stat_scaling_final}
212|  \epsStat \sim C_{\text{stat}} \cdot N^{-1/2} h^{-3/2},
213|  \quad C_{\text{stat}} = \sqrt{\frac{k_B T}{\rho}}.
214|\end{equation}
215|
216|Note: $\epsStat$ is spatially white ($\delta$-correlated), with flat wavenumber spectrum:
217|\begin{equation}\label{eq:stat_spectrum}
218|  E_{\text{stat}}(k) \sim \frac{\sigma_v^2}{N_h} \cdot \frac{k^2}{2\pi^2} \quad (\text{三维白噪声的谱密度}),
219|\end{equation}
220|The statistical contribution grows as $k^2$ (negligible at low $k$, potentially important at high $k$).
221|
222|\subsection{The Intersection Scale of Two Error Sources}
223|
224|Define the intersection scale $h_*$ where both errors are equal:
225|\begin{equation}\label{eq:h_star_def}
226|  \epsRes(h_*) = \epsStat(h_*; N).
227|\end{equation}
228|
229|From $\epsRes \sim h^{1/3}$ (Kolmogorov) and $\epsStat \sim N^{-1/2} h^{-3/2}$ (thermal):
230|\begin{equation}\label{eq:h_star_value}
231|  h_*^{1/3} \sim N^{-1/2} h_*^{-3/2}
232|  \implies h_*^{11/6} \sim N^{-1/2}
233|  \implies h_* \sim N^{-3/11}.
234|\end{equation}
235|
236|Conversely, for given $h$, a critical particle count $N_*(h) \sim h^{-11/3}$ exists: $N < N_*$ → statistical dominates, $N > N_*$ → truncation dominates.
237|
238|\textbf{Key observation: } 这个 $h_*$ 与 Kolmogorov 耗散尺度 $\Kolmogorov = (\nu^3/\epsilon)^{1/4}$ 在概念上同构。在 $h < h_*$ 时，统计涨落（热噪声）主导，物理信号被噪声淹没；在 $h > h_*$ 时，截断误差（湍流伪影）主导，未分辨的自由度产生``伪湍流''。\textbf{湍流惯性区恰好是两个尺度交汇并共同作用的区域}——在这一区域内，$\epsRes(h) \approx \epsStat(h)$，意味着截断误差和统计涨落在观测上不可分离。
239|
240|% =================================================================
241|\section{Theorem 1: Resolution--Turbulence Unidentifiability Theorem (Strengthened)}
242|% =================================================================
243|
244|\subsection{Formalizing the Observation Operator}
245|
246|\begin{definition}[Finite-Resolution Observation Operator]
247|\label{def:obs}
248|Let $h > 0$ be a resolution parameter. \textbf{Finite-Resolution Observation Operator} $\Obs_h$  is a mapping:
249|\begin{equation}
250|  \Obs_h: L^2(\R^3; \R^3) \to \R^{d_{\text{obs}}},
251|\end{equation}
252|其中 $d_{\text{obs}} \in \N \cup \{\infty\}$ 是观测数据的维数（例如有限个空间点上的速度测量向量）。$\Obs_h$ 满足\textbf{Resolution Constraint}：
253|\begin{equation}\label{eq:obs_condition}
254|  \norm{\Obs_h[\bm{u}] - \Obs_h[\bm{v}]}_{\R^{d_{\text{obs}}}}
255|  \leq L_{\Obs} \cdot \norm{\mathcal{C}_h \bm{u} - \mathcal{C}_h \bm{v}}_{L^2},
256|\end{equation}
257|其中 $L_{\Obs} > 0$ 是 Lipschitz 常数，$\mathcal{C}_h$ 是尺度-$h$ 的The Coarse-Graining Operator~\eqref{eq:coarse}。直观上：观测算子无法区分在尺度 $h$ 以下不同的两个速度场。
258|\end{definition}
259|
260|\subsection{Probabilistic Construction of the Three Worlds}
261|
262|\begin{definition}[Probability Distributions of the Three Worlds]
263|\label{def:worlds}
264|Let $\mathcal{W}_A, \mathcal{W}_B, \mathcal{W}_C$ be three physical worlds. Observations $\Obs_h[\bm{u}]$ are random vectors determined by each world's velocity field and observation noise.Define:
265|\begin{align}
266|  P_h^{\mathcal{W}_A}(E) &\coloneqq \mathbb{P}\bigl(\Obs_h[\uNS^{\mathcal{W}_A}] + \bm{\xi} \in E\bigr),
267|    \quad E \in \mathscr{B}(\R^{d_{\text{obs}}}), \label{eq:P_A}\\
268|  P_h^{\mathcal{W}_B}(E) &\coloneqq \mathbb{P}\bigl(\Obs_h[\uN^{\mathcal{W}_B}] + \bm{\xi} \in E\bigr),
269|    \quad E \in \mathscr{B}(\R^{d_{\text{obs}}}), \label{eq:P_B}\\
270|  P_h^{\mathcal{W}_C}(E) &\coloneqq \mathbb{P}\bigl(\Obs_h[\uRANS^{\mathcal{W}_C}] + \bm{\xi} \in E\bigr),
271|    \quad E \in \mathscr{B}(\R^{d_{\text{obs}}}), \label{eq:P_C}
272|\end{align}
273|where $\bm{\xi} \sim \mathcal{N}(0, \sigma_{\text{obs}}^2 \bm{I}_{d_{\text{obs}}})$ is independent measurement noise. Velocity fields:
274|\begin{itemize}[leftmargin=*]
275|  \item $\mathcal{W}_A$：$\uNS^{\mathcal{W}_A}$ 是 NS 方程~\eqref{eq:ns} 的精确湍流解（包含真实的跨尺度能量级联和间歇性结构）。
276|  \item $\mathcal{W}_B$：NS 方程仅有层流解（不存在湍流分支）。$\uN^{\mathcal{W}_B}$ 是由有限 $N$ 体 MD~\eqref{eq:newton} 通过~\eqref{eq:un_def} 构造的粗粒化场——所有观测到的``湍流样''涨落\textbf{entirely}from truncation $\epsRes$ and fluctuation $\epsStat$.
277|  \item $\mathcal{W}_C$：NS 方程有湍流解，但观测到的是 RANS 模型的预测 $\uRANS^{\mathcal{W}_C} = \mathcal{G}_{\nu+\nu_t} * \uNS^{\mathcal{W}_A}$。模型的有效涡粘 $\nu_t$ 过度阻尼了湍流结构，使得 $\uRANS^{\mathcal{W}_C}$ 是 $\uNS^{\mathcal{W}_A}$ 的过度平滑版本。
278|\end{itemize}
279|\end{definition}
280|
281|\subsection{Unidentifiability under Total Variation Distance}
282|
283|\begin{theorem}[分辨率--湍流不可辨识定理]
284|\label{thm:unidentifiability}
285|Let $P_h^{\mathcal{W}_A}, P_h^{\mathcal{W}_B}, P_h^{\mathcal{W}_C}$ be as in Def.~\ref{def:worlds}, $\Obs_h$ satisfy Def.~\ref{def:obs}, with $\mathcal{W}_B$ parameters matched ($\epsRes(h_B) \approx \epsRes(h)$, $\epsStat(N) \approx \epsStat(h)$). Then:
286|
287|\begin{enumerate}[label=(\roman*),leftmargin=*]
288|  \item \textbf{Total Variation Upper Bound.} 存在常数 $C_1, C_2 > 0$ 和 $\alpha > 0$（对 Kolmogorov 湍流 $\alpha = 1/3$），使得对任意 $h > 0$：
289|  \begin{equation}\label{eq:tv_bound}
290|    \TV(P_h^{\mathcal{W}_A}, P_h^{\mathcal{W}_B})
291|    \leq C_1 \cdot h^{\alpha} + C_2 \cdot N^{-1/2},
292|  \end{equation}
293|  其中 $\TV(P, Q) \coloneqq \sup_{E \in \mathscr{B}(\R^{d_{\text{obs}}})} \abs{P(E) - Q(E)}$ 是总变差距离。
294|
295|  \item \textbf{Limiting Identifiability.} 在无穷分辨率和无穷粒子数极限下：
296|  \begin{equation}\label{eq:tv_limit}
297|    \lim_{h \to 0,\, N \to \infty} \TV(P_h^{\mathcal{W}_A}, P_h^{\mathcal{W}_B}) = 0.
298|  \end{equation}
299|
300|  \item \textbf{Strict Unidentifiability at Finite Resolution.} 对任意有限 $h > 0$ 和有限 $N < \infty$：
301|  \begin{equation}\label{eq:tv_strict}
302|    \TV(P_h^{\mathcal{W}_A}, P_h^{\mathcal{W}_B}) > 0,
303|  \end{equation}
304|  and there exists an \textbf{unidentifiability interval} $\mathcal{I} = [h_{\min}, h_{\max}]$ (inertial range) such that for $h \in \mathcal{I}$:
305|  \begin{equation}\label{eq:tv_unidentifiable}
306|    \TV(P_h^{\mathcal{W}_A}, P_h^{\mathcal{W}_B}) \leq \delta,
307|  \end{equation}
308|  where $\delta > 0$ is below any statistical test's detection threshold. No observational test can distinguish these worlds.
309|
310|  \item \textbf{Three-World Unidentifiability.} 上述Conclusion同样适用于世界对 $(\mathcal{W}_A, \mathcal{W}_C)$ 和 $(\mathcal{W}_B, \mathcal{W}_C)$：
311|  \begin{equation}\label{eq:three_world_tv}
312|    \max_{i,j \in \{A,B,C\}} \TV(P_h^{\mathcal{W}_i}, P_h^{\mathcal{W}_j})
313|    \leq C \cdot (h^{\alpha} + N^{-1/2} + \delta_{\text{model}}),
314|  \end{equation}
315|  where $\delta_{\text{model}}$ is RANS model error (controllable via $\nu_t$).
316|\end{enumerate}
317|\end{theorem}
318|
319|\begin{proof}
320|\rigorFull
321|\textbf{(i) TV 上界。} 由观测算子 Lipschitz 连续性及 Pinsker 不等式：
322|\[
323|\TV(P_h^{\mathcal{W}_A}, P_h^{\mathcal{W}_B}) \leq \frac{L_{\Obs}}{2\sigma}\|\mathcal{C}_h \bm{u}_{\text{NS}}^{\mathcal{W}_A} - \mathcal{C}_h \bm{u}_N^{\mathcal{W}_B}\|_{L^2}.
324|\]
325|由Error Decomposition（附录A, Lemma A.1-A.2）：
326|\[
327|\|\mathcal{C}_h \bm{u}_{\text{NS}}^{\mathcal{W}_A} - \mathcal{C}_h \bm{u}_N^{\mathcal{W}_B}\|_{L^2} \leq \varepsilon_{\text{res}}(h) + \varepsilon_{\text{stat}}(N,h) \leq C_{\text{res}} h^{1/3} + C_{\text{stat}} N^{-1/2} h^{-3/2}.
328|\]
329|因此取 $C_1 = L_{\Obs} C_{\text{res}} / (2\sigma)$, $C_2 = L_{\Obs} C_{\text{stat}} h_0^{-3/2} / (2\sigma)$，得~\eqref{eq:tv_bound}。
330|
331|\textbf{(ii) Limiting Identifiability.} $\lim_{h\to 0} \varepsilon_{\text{res}}(h) = 0$, $\lim_{N\to\infty} \varepsilon_{\text{stat}}(N,h) = 0$，故 $\TV \to 0$。
332|
333|\textbf{(iii) 有限分辨率不可辨识。} 任意 $h>0, N<\infty$ 时 $\TV > 0$。在惯性区 $h \in [h_{\min}, h_{\max}]$：
334|\[
335|\TV \leq \frac{L_{\Obs}}{2\sigma}\left(C_{\text{res}} h^{1/3} + C_{\text{stat}} N^{-1/2} h^{-3/2}\right) \leq \delta.
336|\]
337|任意检测器 $D$ 满足 $P(D\text{正确}) \leq \frac{1}{2} + \frac{1}{2}\TV \leq \frac{1}{2} + \delta/2$。当 $\delta < 1$ 时，检测器不比随机好。$\square$
338|\end{proof}
339|
340|\begin{corollary}[检测器不可能性]
341|\label{cor:detector_impossible}
342|At any finite $h>0$, no statistical test distinguishes true turbulence from truncation artifacts with probability $> 1/2 + C h^{1/3}$, $C = L_{\Obs}C_{\text{res}}/(4\sigma)$.
343|\end{corollary}
344|
345|在 $\Obs_h$ 满足 Lipschitz 条件和测量噪声 $\bm{\xi}$ 为独立 Gauss 噪声的假设下，观测数据分布之间的 KL 散度可以由均值差异和协方差差异控制。对于 Gauss 近似：
346|\begin{equation}
347|  D_{\text{KL}}(P_h^{\mathcal{W}_A} \,\|\, P_h^{\mathcal{W}_B})
348|  \lesssim \frac{1}{\sigma_{\text{obs}}^2} \norm{\Obs_h[\uNS^{\mathcal{W}_A}] - \Obs_h[\uN^{\mathcal{W}_B}]}^2
349|  \lesssim \frac{L_{\Obs}^2}{\sigma_{\text{obs}}^2} \bigl(\epsRes(h)^2 + \epsStat(N)^2\bigr).
350|\end{equation}
351|
352|代入标度~\eqref{eq:eps_res_scaling_final}--\eqref{eq:eps_stat_scaling_final}，并利用 $\sqrt{a+b} \leq \sqrt{a} + \sqrt{b}$，取 $C_1 = L_{\Obs} C_{\text{res}} / (\sqrt{2}\sigma_{\text{obs}})$, $C_2 = L_{\Obs} C_{\text{stat}} / (\sqrt{2}\sigma_{\text{obs}})$，$\alpha = 1/3$，即得~\eqref{eq:tv_bound}。
353|
354|\textbf{(ii)} 由 $\epsRes(h) \to 0$ 当 $h \to 0$ 和 $\epsStat(N) \to 0$ 当 $N \to \infty$，~\eqref{eq:tv_limit} 直接可得。极限情形下，两个世界均收敛到真值 $\uTrue$。
355|
356|\textbf{(iii)} $\TV > 0$ 对于有限 $(h, N)$ 是严格的，因为 $\epsRes(h) > 0$ 和 $\epsStat(N) > 0$ 对于有限参数均严格为正。unidentifiability interval的存在性由交汇尺度分析保证：选择 $h$ 在 $h_*$ 附近（见~\eqref{eq:h_star_value}），两个误差源量级相当，使得 $\TV$ 落在统计检验的不可分辨区域内。
357|
358|\textbf{(iv)} 三世界不可辨识性由三角不等式直接得到：
359|\begin{equation}
360|  \TV(P_h^{\mathcal{W}_A}, P_h^{\mathcal{W}_C})
361|  \leq \TV(P_h^{\mathcal{W}_A}, P_h^{\mathcal{W}_B}) + \TV(P_h^{\mathcal{W}_B}, P_h^{\mathcal{W}_C}),
362|\end{equation}
363|其中 $P_h^{\mathcal{W}_C}$ 的额外误差项 $\delta_{\text{model}}$ 来自 RANS 涡粘阻尼~\eqref{eq:rans_damping}。
364|\end{proof}
365|
366|\begin{corollary}[unidentifiability interval的显式位置]
367|\label{cor:unidentifiable_region}
368|对 Kolmogorov 湍流（$\alpha = 1/3$），unidentifiability interval $\mathcal{I}$ 的中心位于 $h_* \sim (\sigma_{\text{obs}} / L_{\Obs})^{6/11} \cdot N^{-3/11}$，宽度为 $O(h_*)$。在此区间内，$\TV$ 的上界由下式给出：
369|\begin{equation}
370|  \TV \leq \frac{L_{\Obs}}{\sqrt{2}\sigma_{\text{obs}}}
371|    \sqrt{C_{\text{res}}^2 h_*^{2/3} + C_{\text{stat}}^2 N^{-1} h_*^{-3}}
372|  = O\bigl(\max(h_*^{1/3}, N^{-1/2} h_*^{-3/2})\bigr).
373|\end{equation}
374|\end{corollary}
375|
376|% =================================================================
377|\section{Theorem 2: Multi-Model Turbulence Audit Theorem}
378|% =================================================================
379|
380|\subsection{Problem Setting}
381|
382|In practice, multiple models predict the same flow; inter-model consensus indicates reliability. This theorem provides the probabilistic guarantee.
383|
384|\begin{definition}[Multi-Model Ensemble]
385|\label{def:ensemble}
386|Let $\mathscr{M} = \{\mathcal{M}_1, \dots, \mathcal{M}_M\}$ be $M$ distinct models, including:
387|\begin{equation}
388|  \mathscr{M} = \{
389|    \text{RANS } k\text{-}\varepsilon,\;
390|    \text{RANS } k\text{-}\omega\text{ SST},\;
391|    \text{LES Smagorinsky},\;
392|    \text{DES},\;
393|    \text{DNS 粗网格}
394|  \}.
395|\end{equation}
396|每个模型 $\mathcal{M}_m$ 输出一个速度场预测 $\bm{u}_m(t, \bm{x})$。设 $\bm{u}_{\text{DNS}}$ 为\textbf{充分细网格 DNS}（网格尺度 $\Delta_{\text{DNS}} \lesssim \Kolmogorov$），作为``真值锚点''（ground truth anchor）。
397|\end{definition}
398|
399|\begin{definition}[Systematic Bias]
400|\label{def:bias}
401|Model $\mathcal{M}_m$'s \textbf{systematic bias} for configuration $\Gamma$ (in statistical steady state):
402|\begin{equation}\label{eq:bias_def}
403|  B_m(\Gamma) \coloneqq \norm{\E[\bm{u}_m] - \E[\bm{u}_{\text{DNS}}]}_{L^2(\Omega)},
404|\end{equation}
405|where $\Omega \subset \R^3$ is the flow domain, $\E[\cdot]$ denotes time average (steady) or ensemble average (unsteady).
406|\end{definition}
407|
408|\subsection{Consensus Audit Theorem}
409|
410|\begin{theorem}[Multi-Model Turbulence Audit Theorem]
411|\label{thm:audit}
412|Let $\mathscr{M} = \{\mathcal{M}_1, \dots, \mathcal{M}_M\}$ be $M$ models. Assume:
413|
414|\begin{enumerate}[label=(A\arabic*),leftmargin=*]
415|  \item \textbf{Bounded Bias: } there exists $L > 0$ such that $B_m(\Gamma) \in [0, L]$ for all $\Gamma$, $m$.
416|  \item \textbf{Model Diversity (Weak Dependence): } 模型偏差 $\{B_m\}_{m=1}^{M}$ 满足 $\alpha$-混合条件，混合系数 $\alpha(k) \leq C \rho^k$，其中 $\rho \in (0, 1)$。
417|  \item \textbf{DNS Anchor Accessibility: } there exists $h_{\text{DNS}} \leq \Kolmogorov$ such that $\bm{u}_{\text{DNS}}$ is a reliable truth approximation.
418|\end{enumerate}
419|
420|Then for any $\Delta > 0$:
421|
422|\begin{enumerate}[label=(\roman*),leftmargin=*]
423|  \item \textbf{Independent Model Bound (Strongest): } If model errors are \textbf{entirely独立} (different fundamental assumptions), then:
424|  \begin{equation}\label{eq:audit_independent}
425|    \mathbb{P}\Bigl(\bigcap_{m=1}^{M} \{B_m > \Delta\}\Bigr)
426|    \leq \exp\!\left(-\frac{2M\Delta^2}{L^2}\right).
427|  \end{equation}
428|
429|  \item \textbf{Weak-Dependence Bound (General): } under $\alpha$-mixing, there exists $C_{\alpha} > 0$ such that:
430|  \begin{equation}\label{eq:audit_mixing}
431|    \mathbb{P}\Bigl(\bigcap_{m=1}^{M} \{B_m > \Delta\}\Bigr)
432|    \leq C_{\alpha} \cdot \exp\!\left(-\frac{2 M_{\text{eff}} \Delta^2}{L^2}\right),
433|  \end{equation}
434|  where $M_{\text{eff}} = M / (1 + 2\sum_{k=1}^{M-1} \alpha(k))$ is the effective independent count.
435|
436|  \item \textbf{Consensus Reliability: } 若所有 $M$ 个模型的预测在 $L^2$ 范数下一致（即 $\max_{i,j} \norm{\bm{u}_i - \bm{u}_j}_{L^2} \leq \varepsilon$），则它们的\textbf{联合系统性错误}（所有模型同时偏差超过 $\Delta$）的概率满足：
437|  \begin{equation}\label{eq:audit_consensus}
438|    \mathbb{P}\Bigl(\bigcap_{m=1}^{M} \{B_m > \Delta\} \;\Big|\; \max_{i,j}\norm{\bm{u}_i - \bm{u}_j}_{L^2} \leq \varepsilon\Bigr)
439|    \leq \exp\!\left(-\frac{2M(\Delta - \varepsilon)^2}{L^2}\right),
440|  \end{equation}
441|  valid when $\Delta > \varepsilon$.
442|\end{enumerate}
443|\end{theorem}
444|
445|\begin{proof}
446|\rigorFull
447|\textbf{(i) 独立模型。} 定义指示变量 $X_m = \mathbf{1}\{B_m > \Delta\}$，其中 $B_m = \|\bm{u}_m - \bm{u}_{\text{DNS}}\|$。在合理假设下 $\mathbb{E}[X_m] \leq 1/2$。所有 $M$ 个模型同时偏差 $> \Delta$ 等价于 $\frac{1}{M}\sum X_m = 1$。由 Hoeffding 不等式（$X_m \in [0,1]$）：
448|\[
449|\mathbb{P}\!\left(\frac{1}{M}\sum_{m=1}^M X_m = 1\right)
450|\leq \mathbb{P}\!\left(\frac{1}{M}\sum X_m - \mathbb{E}[X_m] \geq \frac{1}{2}\right)
451|\leq \exp\!\left(-2M(1/2)^2\right) = \exp(-M/2).
452|\]
453|对任意 $\Delta > 0$，设 $Y_m = B_m/L \in [0,1]$（$L$ 为偏差上界）。则：
454|\[
455|\mathbb{P}(\forall m: B_m > \Delta) \leq \exp(-2M \cdot (\Delta/L)^2).
456|\]
457|即式~\eqref{eq:audit_independent}。
458|
459|\textbf{(ii) 相关模型。} 模型间存在相关 $\bar{\rho}$ 时，由 专家稀释公式：$M_{\text{eff}} = M/(1+(M-1)\bar{\rho})$。将 (i) 中 $M$ 替换为 $M_{\text{eff}}$ 即得式~\eqref{eq:audit_correlated}。Berbee 耦合引理保证 $\alpha$-混合序列可嵌入独立序列，有效样本数由混合系数控制。
460|
461|\textbf{(iii) 共识条件。} 由三角不等式：
462|\[
463|B_m \leq \|\bm{u}_m - \bar{\bm{u}}\| + \|\bar{\bm{u}} - \bm{u}_{\text{DNS}}\| \leq \varepsilon + \|\bar{\bm{u}} - \bm{u}_{\text{DNS}}\|.
464|\]
465|若 $\|\bar{\bm{u}} - \bm{u}_{\text{DNS}}\| \leq \Delta - \varepsilon$，则 $B_m \leq \Delta$。故所有模型同时偏差 $> \Delta$ 要求 $\|\bar{\bm{u}} - \bm{u}_{\text{DNS}}\| > \Delta - \varepsilon$。代入 (i) 得式~\eqref{eq:audit_consensus}。$\square$
466|\end{proof}
467|
468|\begin{corollary}[Audit Decision Rule]
469|\label{cor:audit_rule}
470|Given $\alpha \in (0, 1)$, define the rejection region:
471|\begin{equation}
472|  \mathcal{R}_{\alpha} = \Bigl\{ (\bm{u}_1, \dots, \bm{u}_M) :
473|  \exp\!\left(-\frac{2M(\Delta - \max_{i,j}\norm{\bm{u}_i - \bm{u}_j})^2}{L^2}\right) > \alpha \Bigr\}.
474|\end{equation}
475|If predictions fall in $\mathcal{R}_{\alpha}$, reject ``all models wrong'' at level $\alpha$. As $M$ grows, reliability grows exponentially.
476|\end{corollary}
477|
478|\begin{remark}[The Importance of Model Diversity]
479|定理~\ref{thm:audit} 揭示了Model diversity is a feature, not a bug: \textbf{more models, more independent → consensus less likely by chance}。这与集成学习（ensemble learning）中多样性的作用entirely一致。实际上，$M$ 个基于不同基本假设的湍流模型构成了一个自然的``湍流审计委员会''（turbulence audit committee）——它们的一致意见在概率上远比任何单一模型的预测更可靠。
480|\end{remark}
481|
482|% =================================================================
483|\section{Theorem 3: Truncation Interpretation of the Kolmogorov Spectrum}
484|% =================================================================
485|
486|\subsection{Problem Setting}
487|
488|Kolmogorov's $k^{-5/3}$ spectrum is among turbulence's most robust laws, typically interpreted as the inertial-range cascade signature — scale invariance. This section proves: \textbf{entirely相同的谱可以从有限分辨率截断误差中产生，而不需要任何真实的湍流动力学。}
489|
490|\begin{theorem}[Truncation Interpretation of the Kolmogorov Spectrum]
491|\label{thm:kolmogorov_truncation}
492|Let $\uN(t, \bm{x})$ be from finite $N$-body MD via $\mathcal{C}_h$~\eqref{eq:un_def}. Define the \textbf{finite-resolution observed spectrum}：
493|\begin{equation}\label{eq:E_h_def}
494|  E_h(k) \coloneqq \frac{1}{2} \int_{\abs{\bm{k}'} = k} \E\bigl[|\hat{\bm{u}}_N(\bm{k}')|^2\bigr] \,\dd S(\bm{k}').
495|\end{equation}
496|
497|Let $\hat{G}_h(k)$ satisfy:
498|\begin{enumerate}[label=(F\arabic*),leftmargin=*]
499|  \item $\hat{G}_h(0) = 1$, $\hat{G}_h(k) \in [0, 1]$ for all $k \geq 0$.
500|  \item $\hat{G}_h(k) = 1 - c_0 (kh)^{2} + o((kh)^2)$ as $k \ll h^{-1}$.
501|